A database engine must answer every statement or reject it with an error. Mature
engines check this with differential fuzzers that compare their answers against
a reference engine. Such a fuzzer cannot test a function it never generates, nor
judge one its reference engine lacks. Turso, a SQLite-compatible engine written
in Rust, runs such a fuzzer in continuous integration, yet the single statement
\sql{select LPAD('a', 9223372036854775807, 'b')} terminated the whole process.
We present \tool. It takes the shape of bug reports filed against a different
engine, DuckDB, refills their literals with boundary values, and runs them
against Turso together with a typed SQL generator. Its only oracle needs no
reference engine: one statement must not end the process. \tool found three
previously unknown defects in Turso's \sql{lpad}, \sql{rpad} and \sql{repeat}
functions. Each aborts the engine from one statement, and \tool regenerates each
from a recorded seed. Our patch for them exposed a fourth defect in the same
length check, which \tool did not find. Our two patches, which fix all four,
were merged upstream. None of Turso's own fuzzers can express any of the four,
since the three functions are absent from all of them. At an equal budget of
20{,}000 cases, mined reports and typed generation each reached both root causes
of the three defects, and drawing from both sources found the first defect at
case 101. The work taught us that reading an incumbent fuzzer's source is a cheap
and reliable way to decide where to test next.
